\documentclass[10pt,a4paper]{amsart}
\usepackage[margin=19mm]{geometry}
\usepackage{amsmath,amssymb,mathtools}
\usepackage[scr=rsfs,cal=euler]{mathalfa}
\usepackage{xfrac}
\usepackage[unicode,hidelinks,bookmarks=false]{hyperref}
\newcommand{\ie}{i.e.\ }
\newcommand{\cf}{cf.\ }
\newcommand{\resp}{respectively\ }
\newcommand{\Subsection}{\subsection}
\providecommand{\nobreakdash}{\nobreak\hbox{-}}
\setlength{\emergencystretch}{2em}
\multlinegap0pt
\usepackage{amssymb}
\usepackage{mathtools}
\usepackage{algorithm}\usepackage{algpseudocode}
\usepackage{bm}
\newtheorem{lem}{Lemma}
\newtheorem{prop}{Proposition}
\newcommand{\Id}{\mathrm{Id}}
\renewcommand{\Im}{\mathfrak{Im}} % imaginary part
\newcommand{\Vh}{\mathrm{V}_{h}}
\newcommand{\bp}{\boldsymbol{p}}
\newcommand{\bq}{\boldsymbol{q}}
\newcommand{\bu}{\boldsymbol{u}}
\newcommand{\bv}{\boldsymbol{v}}
\newcommand{\bw}{\boldsymbol{w}}
\newcommand{\mA}{\mathrm{A}}
\newcommand{\mB}{\mathrm{B}}
\newcommand{\mH}{\mathrm{H}}
\newcommand{\mL}{\mathrm{L}}
\newcommand{\mR}{\mathrm{R}}
\newcommand{\mS}{\mathrm{S}}
\newcommand{\mT}{\mathrm{T}}
\newcommand{\mbV}{\mathbb{V}}
\newcommand{\mbX}{\mathbb{X}}
\newcommand{\mrm}{\mathrm}
\title{\bf Discrete FEM-BEM coupling with the\\ Generalized Optimized Schwarz Method}
\author{Antonin Boisneault, Marcella Bonazzoli, Xavier Claeys, Pierre Marchand}
\date{Source: arXiv:2601.16817v1 (CC BY 4.0)}
\begin{document}
\maketitle

\noindent\emph{Source and licence.} \bf Discrete FEM-BEM coupling with the\\ Generalized Optimized Schwarz Method. Version arXiv:2601.16817v1 (CC BY 4.0), distributed by arXiv under the Creative Commons Attribution 4.0 International licence, http://creativecommons.org/licenses/by/4.0/, as recorded in the arXiv licence metadata for this exact version and shown on the abstract page https://arxiv.org/abs/2601.16817v1. Full licence notice: \texttt{licenses/arxiv-2601.16817-CC-BY-4.0.txt}.\par\smallskip

\noindent\textit{Editorial scope.} Counted unit: the article's prop IsomorphismKernels2 (source0.tex lines 1633--1637). The article's complete original proof of it is delivered with it (source0.tex lines 1638--1690, one \texttt{proof} environment of 53 lines). No mathematics is written, repaired or re-derived: the statement and the proof are the article's own lines, in the article's own notation, and no retained line is substituted or re-flowed.  Retained uncounted as the source-local context the proof needs: lines 424--426; lines 642--648; lines 1130--1135; lines 1144--1149; lines 1199--1202; lines 1285--1294; lines 1417--1420; lines 1469--1474; lines 1596--1602; lines 1613--1615.  Nothing was omitted from any retained span, so the package carries no omission record.  No retained line contains a citation, so the wrapper carries no bibliography and no bibliography entry.  No retained line contains a figure environment, an image-inclusion command or drawing code, so no image is delivered and the package declares no asset dependency.  Editorial formatting only, in addition: an AMS \texttt{amsart} preamble that restates the article's own declarations and macros used by the retained lines, namely the environments \texttt{lem}, \texttt{prop} on separate counters, together with the standard packages those lines need.  The retained environments print in this document as Lemma 1, Proposition 1 in order of appearance, whereas the article prints the same items with the numbering of its own full document.  The complete native source of the article as distributed in the arXiv e-print for this exact version is bundled unchanged as \texttt{sources/193263/source0.tex}.
\begin{align}\label{eq:polarOfImage}
  \mrm{Im}(\mL)^{\circ} = \mrm{Ker} (\mL^*).
\end{align}

\begin{equation}\label{Assumption2}
  \begin{aligned}
    & \textbf{Assumption:}\\
    & \Im\{\langle \mathcal{A}_{\Gamma}(v,q),(\overline{v},\overline{q})\rangle\}\leq 0\\
    & \forall v\in\mH^{+\frac{1}{2}}(\Gamma),\;\;\forall q\in\mH^{-\frac{1}{2}}(\Gamma).
  \end{aligned}
\end{equation}

\begin{equation}\label{DefinitionSingleTraceSpace}
  \begin{aligned}
    & \mbX_h(\Omega)\coloneq \mathcal{R}(\Vh(\Omega,\Gamma)),\\
    & \mbX_h(\Sigma)\coloneq \mR(\Vh(\Sigma)).
  \end{aligned}
\end{equation}

\begin{lem}\label{ReformulationTransmission}\quad\\[-15pt]
  \begin{itemize}
  \item[i)]  $\mbX_h(\Omega) = \mB^{-1}(\mbX_h(\Sigma))$,\\[-20pt]
  \item[ii)] $\mbX_h(\Omega)^{\circ} =\mB^{*}(\mbX_h(\Sigma)^{\circ})$.
  \end{itemize}
\end{lem}

  \begin{equation}\label{CaracAvecEchangeOperator}
    (\bv,\bp)\in \mathscr{X}_h(\Sigma) \coloneq \mbX_h(\Sigma)\times\mbX_h(\Sigma)^{\circ}\quad
    \iff\quad -\bp+i\mT(\bv) = \Pi(\bp+i\mT(\bv)).
  \end{equation}

\begin{lem}\label{ElementaryPropScattering}\quad\\
  If either \(\mA-i\mB^*\mT\mB\) is invertible, or Assumption~\eqref{Assumption2} holds, then\\[-5pt]
  \begin{equation*}
    \hspace{-8.65cm}
    \begin{aligned}
     & \textit{i)}  & \mrm{Ker}(\mA-i\mB^*\mT\mB) & = \mrm{Ker}(\mA)\cap\mrm{Ker}(\mB)\\
     & \textit{ii)} & \mrm{Im}(\mA-i\mB^*\mT\mB)  & = \mrm{Im}(\mA)+\mrm{Im}(\mB^*)\\[5pt]
    \end{aligned}
  \end{equation*}
\end{lem}

  \begin{equation}\label{ScatteringOperator}
    \bp+i\mT\bv = \mS(\bp-i\mT\bv),\quad \forall (\bv,\bp)\in
    \mathscr{C}_h(\mA).
  \end{equation}

\begin{lem}\label{ScatteringOperatorChara}\quad\\
  Assume one of the hypotheses of Lemma~\ref{ElementaryPropScattering}. For any \((\bv,\bp)\in \mathscr{H}_h(\Sigma)\)
  \begin{align*}
    (\bv,\bp)\in \mathscr{C}_h(\mA) \Longleftrightarrow \bp+i\mT\bv = \mS(\bp-i\mT\bv). 
  \end{align*}
\end{lem}

\begin{prop}\label{IsomorphismKernels1}\quad\\
  Assume one of the hypotheses of Lemma~\ref{ElementaryPropScattering}. Define the operator $\Phi\colon 
  \mrm{Ker}(\mA_{\Omega\times\Gamma})\to \mbV_h(\Sigma)^*$ by the formula
  $\Phi(\bu) \coloneq (\mB^{\dagger})^*\cdot(\mA-i\mB^*\mT\mB)\cdot\mathcal{R}(\bu)$.
  Then $\bu\in \mrm{Ker}(\Phi)\iff\mathcal{R}(\bu)\in \mrm{Ker}(\mA)\cap \mrm{Ker}(\mB)$, 
  and $\mrm{dim}\,\mrm{Ker}(\Phi) = \mrm{dim}(\mrm{Ker}(\mA)\cap \mrm{Ker}(\mB))$.
\end{prop}

\begin{equation}\label{IdentitePhiIntermediaire}
  \Phi(\bu) = \bp-i\mT(\bv).
\end{equation}

\begin{prop}\label{IsomorphismKernels2}\quad\\
  Assume one of the hypotheses of Lemma~\ref{ElementaryPropScattering}. 
  Let $\Phi$ be the operator defined in Proposition~\ref{IsomorphismKernels1}. 
  Then $\mrm{Im}(\Phi) = \Phi(\,\mrm{Ker}(\mA_{\Omega\times\Gamma})\,) = \mrm{Ker}(\Id+\Pi\mS)$.
\end{prop}

\begin{proof}
Recall the description of the map $\Phi$ from the beginning of the
proof of Proposition~\ref{IsomorphismKernels1} and~\eqref{IdentitePhiIntermediaire}. What precedes shows that, if $\bu\in
\mrm{Ker}(\mA_{\Omega\times\Gamma})$, then $\Phi(\bu) = \bp-i\mT(\bv)$
for some $(\bv,\bp)\in \mathscr{C}_h(\mA)\cap\mathscr{X}_h(\Sigma)$.
Applying the definition of the scattering operator $\mS$, see~\eqref{ScatteringOperator}, $\mS\cdot\Phi(\bu) =
\bp+i\mT(\bv)$. Then, applying the characterization of $\mathscr{X}_h(\Sigma)$ with the exchange operator
$\Pi$, see~\eqref{CaracAvecEchangeOperator}, we obtain
$\Pi\cdot\mS\cdot\Phi(\bu) = \Pi(\bp+i\mT(\bv)) = -\bp+i\mT(\bv) =
-\Phi(\bu)$. Hence, $(\Id + \Pi\mS)\Phi(\bu) = 0$ for all \(\bu\in
\mrm{Ker}(\mA_{\Omega\times\Gamma})\), which rewrites
\begin{equation*}
  \Phi(\,\mrm{Ker}(\mA_{\Omega\times\Gamma})\,)\subset \mrm{Ker}(\Id + \Pi\mS).
\end{equation*}

There only remains to show that
$\Phi(\,\mrm{Ker}(\mA_{\Omega\times\Gamma})\,) \supset \mrm{Ker}(\Id +
\Pi\mS)$. Pick an arbitrary $\bq_-\in \mrm{Ker}(\Id + \Pi\mS)$ and
set $\bq_+ = \mS(\bq_-)$. Since $(\Id + \Pi\mS)\bq_- = 0$ we
conclude that $\bq_-+\Pi(\bq_+) = 0$. Next define $(\bv,\bp)\in
\mbV_h(\Sigma)\times \mbV_h(\Sigma)^*$ by
\begin{equation*}
  \left\{
  \begin{array}{l}
    \bp-i\mT\bv = \bq_-\\
    \bp+i\mT\bv = \bq_+
  \end{array}
  \right.
  \quad\iff\quad
  \left\{
  \begin{array}{l}
    \bv = \mT^{-1}(\bq_+ -\bq_-)/(2i)\\
    \bp = (\bq_+ + \bq_-)/2
  \end{array}
  \right.
  .
\end{equation*}
According to the equations satisfied by $(\bq_-,\bq_+)$, we
have $\bp+i\mT\bv = \mS(\bp-i\mT\bv)$ and $-\bp+i\mT\bv =
\Pi(\bp+i\mT\bv)$, which rewrites
$(\bv,\bp)\in\mathscr{C}_h(\mA)\cap\mathscr{X}_h(\Sigma)$ due to~\eqref{CaracAvecEchangeOperator} and Lemma~\ref{ScatteringOperatorChara}.
Since $(\bv,\bp)\in\mathscr{C}_h(\mA)$, there exists $\bw\in
\mbV_h(\Omega)$ such that $\mB(\bw) = \bv$ and $\mA\bw =
\mB^*(\bp)$. On the other hand, as $\bv\in\mbX_h(\Sigma)$, we derive
$\bw\in \mB^{-1}(\mbX_h(\Sigma)) = \mbX_h(\Omega) =
\mathcal{R}(\Vh(\Omega,\Gamma))$, see \textit{i)} of Lemma~\ref{ReformulationTransmission} and~\eqref{DefinitionSingleTraceSpace}. So $\bw = \mathcal{R}(\bu)$ for
some $\bu\in \Vh(\Omega,\Gamma)$, and
$\mA_{\Omega\times\Gamma}(\bu) = \mathcal{R}^*\mA\mathcal{R}(\bu) =
\mathcal{R}^*\mA\bw = \mathcal{R}^*\mB^*\bp = \mathcal{B}^*R^*\bp = 0$ since $\bp\in
\mbX_h(\Sigma)^{\circ} = (\mR(\Vh(\Sigma)))^\circ = \mrm{Ker}\,\mR^*$ by~\eqref{eq:polarOfImage}. As a consequence, $\bu \in
\mrm{Ker}(\mA_{\Omega\times\Gamma})$ and $\Phi(\bu) = \bp-i\mT(\bv) =
\bq_-$. 
\end{proof}

\end{document}
